Ce projet AniHub a consisté à concevoir et développer une plateforme web à architecture microservices dédiée aux animaux domestiques en Tunisie. Il vise à centraliser plusieurs services essentiels dans un seul système unifié.

La plateforme permet la gestion d’annonces d’achat et de vente, les demandes d’adoption, la mise en relation pour l’accouplement, ainsi que divers services liés aux animaux (toilettage, soins, garde). Elle intègre également un module de déclaration de perte d’animaux afin de faciliter la recherche et d’encourager l’entraide entre utilisateurs.

Grâce à l’approche microservices, l’application bénéficie d’une meilleure modularité, scalabilité et maintenabilité, tout en garantissant une évolution progressive des fonctionnalités.

AniHub représente ainsi une solution complète répondant à un besoin réel. \newline
Dans cette optique, plusieurs perspectives d’amélioration sont envisagées :

- \textbf{Amélioration de l’expérience utilisateur (UX/UI)} : il s’agit d’optimiser l’interface de la plateforme afin de la rendre plus intuitive, fluide et moderne. Cela permettrait de simplifier la navigation et d’améliorer l’expérience globale des utilisateurs.

- \textbf{Développement d’une application mobile AniHub} : cette perspective vise à étendre la plateforme vers le mobile afin de permettre un accès plus rapide et plus pratique depuis les smartphones, améliorant ainsi l’accessibilité et l’engagement des utilisateurs.

- \textbf{Mise en place de deux plans d’utilisation (Free et Premium)} : l’objectif est d’introduire un modèle économique différencié, où le plan gratuit offre les fonctionnalités de base tandis que le plan premium propose des services avancés et des options supplémentaires.

- \textbf{Intégration d’un module d’intelligence artificielle (IA)} : cette amélioration permettra d’ajouter des fonctionnalités intelligentes telles que la recommandation personnalisée des annonces et la détection automatique des annonces frauduleuses ou inappropriées, afin d’améliorer la qualité et la sécurité de la plateforme.